\chapter{Future Work}

Although the proposed workflow automation platform successfully demonstrates AI-assisted workflow generation and execution, several opportunities remain for future enhancement. One promising direction is expanding the library of supported nodes to cover a wider range of third-party services, cloud platforms, and IoT devices. Supporting standards such as OPC UA, Modbus, and additional industrial communication protocols would make the platform more suitable for real-world automation and smart manufacturing applications.

Another important area for improvement is the AI workflow generation pipeline. While the current implementation relies on multiple specialized agents to interpret user requests and construct workflows, future versions could incorporate reinforcement learning from user feedback and execution history to improve planning accuracy. Integrating retrieval-augmented generation (RAG) with external documentation and node descriptions would also enable the agent to generate more reliable and context-aware workflows.

The execution engine can also be extended to support distributed execution across multiple worker instances and machines. Such an enhancement would improve scalability, fault tolerance, and throughput for computationally intensive workflows. Additional monitoring capabilities, including real-time execution visualization, performance analytics, and automatic failure recovery, would further improve system reliability and simplify workflow debugging.

Finally, future work may focus on introducing advanced workflow optimization techniques. These include automatic parallelization of independent tasks, resource-aware scheduling, workflow versioning, and collaborative editing. Combined with stronger security features such as fine-grained access control, audit logging, and role-based permissions, these improvements would transform the platform from a research prototype into a production-ready workflow automation solution suitable for enterprise-scale deployments.